\Needspace{15\baselineskip}
\section{The Eucharistic Prayer: One Address, Many Movements}

The preparation of altar and gifts ends in a prayer over what has been set apart. The Eucharistic Prayer then changes the scale of the action. It is not a priestly devotion performed while the congregation waits for Communion, nor a narration addressed to the congregation. It is the Church's ``prayer of thanksgiving and sanctification,'' voiced by the ordained Priest to the Father through Christ in the Holy Spirit, while the baptized assembly joins Christ in confessing God's works and offering the sacrifice (GIRM 78). The dialogue through the Great Amen is one address. The following movements distinguish its actions without breaking that unity.

\ordermovement{\latin{Dialogus praefationis} and the Preface}{Hearts Raised into Particular Thanksgiving}{Fixed dialogue followed by the proper or lawfully selected Preface; some Eucharistic Prayers require their own Preface}{GIRM 78--79a, 147--148, 364--365; \latin{Ordo Missae}}

Standing, the Priest and people exchange the Preface dialogue. The greeting renews their liturgical relation; the summons concerning the heart names prayer's direction; the final exchange acknowledges thanksgiving as right. The Priest then continues with hands extended. A proper, seasonal, common, or prayer-specific Preface gives the thanksgiving its received content. The dialogue is therefore stable while the cause for praise is often variable. Selection is not thematic decoration: the calendar, Mass formulary, and Eucharistic Prayer determine which Preface may govern the prayer.

\paragraph{What is raised?}
Lamentations~3:40--42 places raised heart and hands after examination and return; Colossians~3:1--4 directs those raised with Christ toward the life hidden in him and awaiting manifestation. The liturgical ascent is accordingly neither contempt for creation nor escape from history. Bread, wine, bodies, grief, labor, and the day's proper mystery remain present. They are reoriented to the Father in Christ. The imperative is possible because grace has first raised the worshippers with the Son.

Cyprian supplies a secure mid-third-century witness to the dialogue and gives it an ethical edge: when the people answer that their hearts are directed to the Lord, their attention must not contradict their speech (\emph{De dominica oratione} 31). Augustine returns to the exchange in preaching to the newly baptized. Hearts are above, he argues, not as a claim of spiritual achievement but because God lifts them; thanksgiving follows because the ascent itself is gift (Sermon 227). These African witnesses do not prove the date of every Roman Preface. They do show that the dialogue had already become a school of collected desire, received grace, and truthful assent.

\paragraph{Why does thanksgiving need a proper cause?}
The Preface is not merely an ornate runway leading to the Sanctus. Its recurring judgment that thanksgiving is fitting opens into a particular reason: creation, redemption, a season, a mystery of Christ, the grace manifested in a saint, or the thematic economy integral to a particular Eucharistic Prayer. Cyril of Jerusalem describes lifted hearts followed by thanksgiving for God's works and the angelic hymn (\emph{Mystagogical Catechesis} V.4--6). He is not the author of the present Roman corpus. His local witness nevertheless displays an old anaphoral grammar: the Church names divine beneficence before joining the praise of heaven.

Roman Prefaces accumulated, contracted, and were restored across many centuries. Early sacramentaries preserve a wide and uneven treasury; the post-Tridentine Missal carried a smaller set; twentieth-century reforms and the postconciliar Missal recovered and composed many more. A common opening does not make their theology interchangeable. A fixed salvation-history Preface is structurally inseparable from Eucharistic Prayer IV, while Eucharistic Prayer III can receive the thanksgiving proper to the day. Eucharistic Prayer II has an own Preface yet admits qualifying alternatives. The two Reconciliation prayers admit a narrower class of penitential or conversion Prefaces. The four prayers for Various Needs and the children's prayers carry integral Prefaces under their own rules. The heart is lifted into an actual ecclesial mystery, not into a generic religious elevation.

\begin{studybox}{The Preface belongs to the great prayer}
The familiar English label can suggest preliminary material. Liturgically the Preface is the Eucharistic Prayer's thanksgiving, not an introduction outside it. A fixed Preface and a proper Preface differ in selection, but both establish what the prayer is thanking the Father for through Christ.
\end{studybox}

\ordermovement{\latin{Sanctus}}{The Earthly Church Enters the Heavenly Hymn}{Fixed acclamation sung or said by the Priest together with all the people}{GIRM 79b, 148; \latin{Ordo Missae}}

The Preface does not report that angels praise and then leave the assembly as audience. Its conclusion invites earthly voices into their hymn. Priest and people together sing or say the complete acclamation; it is not an ordinary presidential sentence with a short response. Yet it remains inside the one Eucharistic Prayer. Its final cry does not conclude the thanksgiving as though a second prayer must begin afterward.

\paragraph{Holiness, judgment, and mission.}
The first scriptural field is Isaiah~6:1--8. The threefold acclamation surrounds the enthroned Lord; thresholds shake, smoke fills the house, and holiness exposes the prophet's unclean lips. Fire from the altar purifies him, and adoration becomes sending. Revelation~4:1--11 places the same ceaseless praise around the heavenly throne and draws elders and creation into worship. The liturgical Sanctus therefore does more than multiply an adjective. It confesses divine otherness as glory filling creation, judgment revealing sin, mercy cleansing worshippers, and holiness issuing in mission.

Cyril explicitly says that the Church sings the seraphic hymn in order to share the praise of the heavenly hosts (\emph{Mystagogical Catechesis} V.5--6). The fourth-century \emph{Apostolic Constitutions} VIII.12 gives another developed Eastern Eucharistic witness; the early Roman sacramentaries provide the first secure transmission in the Roman books. Such evidence establishes reception in distinct churches, not a single route of borrowing or a named composer of the Roman text. A later \emph{Liber pontificalis} attribution to an early pope cannot bear the weight sometimes placed upon it.

\paragraph{The holy one comes.}
The second scriptural field joins Psalm~118:25--26 to the entry into Jerusalem (Matt.~21:9; Mark~11:9--10; John~12:13). Temple procession, royal advent, and the approach to the Passion converge. The one adored beyond all creaturely measure comes in the Lord's name and gives himself. This conjunction fittingly intensifies expectation before the consecratory center, but it is not a prediction that the consecration has not yet ``started,'' nor does the acclamation effect the sacramental change. It holds transcendence and advent together: the thrice-holy Lord is the crucified and risen king who approaches his Church.

Because the whole people acclaim, musical settings should serve intelligible corporate praise rather than turn the assembly into spectators of a specialist piece. Because the acclamation belongs to the Eucharistic Prayer, replacing it with another hymn changes the received anaphora rather than merely changing music. The Church does not invent a bridge between Preface and post-Sanctus; she receives words in which prophetic temple, Gospel advent, and heavenly liturgy already meet.

\ordermovement{The Presidential Voice and the Assembly's Participation}{One Body Participates through Differentiated Voices}{The Eucharistic Prayer is proclaimed in full by the Priest alone; concelebrants take assigned parts; the assembly joins through faith, silence, posture, and the prescribed acclamations}{GIRM 30--34, 78, 147, 216--236; \emph{Redemptionis Sacramentum} 51--54}

The Priest addresses the Father in the name of the whole holy people. At a concelebration the missal distributes particular passages among ordained concelebrants and supplies the manner of their common utterance. No rubric thereby gives the prayer to the congregation for collective recitation. Conversely, priestly presidency does not make the people external to it. GIRM 78 says that the entire congregation joins itself with Christ in praise and offering; GIRM 147 names reverent silence and the appointed acclamations as modes of that union.

\paragraph{Whose ``we'' is heard?}
Christ is the principal agent of the liturgy and the one High Priest. The ordained Priest acts in Christ's person as Head and voices an ecclesial prayer he did not compose. The baptized exercise their common priesthood by bringing faith, praise, self-offering, intercession, and sacramental reception into Christ's oblation. \emph{Lumen gentium} 10 insists that ministerial and common priesthood differ in essence and not merely degree, while being ordered to one another. The anaphoral plural neither makes the celebrant a dramatic impersonator of a crowd nor collapses distinct offices into a common speaking role.

Justin Martyr describes a president offering thanksgiving and prayers and the people sealing them with Amen (\emph{First Apology} 65, 67). His account predates the fixed Roman texts and does not describe every present posture or acclamation. It nonetheless witnesses a durable polarity: presidential thanksgiving is genuinely the assembly's prayer, and corporate ratification does not require identical utterance. Cyprian's demand that the heart match the Preface response applies throughout the sustained address.

\paragraph{Silence as an ecclesial act.}
The silence of the Eucharistic Prayer differs from private recollection before Mass or contemplation after Communion. It is concentrated listening to a prayer in which the hearer is implicated. The people have brought gifts, answered the dialogue, and sung the Sanctus; they will acclaim the memorial and ratify the doxology. Between those utterances they attend to the Father's praise, the Spirit's invocation, Christ's words, the Church's offering, and intercession for living and dead. ``Active'' participation cannot be measured by the percentage of syllables assigned to each person. Listening can be more demanding than simultaneous speech because it requires the hearer to let another's received words become his own prayer.

The rubrical discipline protects both sides. A Priest may not hand his presidential part to a lay person or invite improvised collective recitation; the faithful are not to be treated as an audience whose only meaningful act begins at Communion. Amplification and music serve audibility and unity rather than theatrical personality. The voice belongs to an office, and the prayer belongs to the Church.

\ordermovement{Epiclesis}{The Father Is Asked to Sanctify Gift and People by the Spirit}{A constitutive anaphoral action expressed according to each Eucharistic Prayer's received wording and order}{GIRM 79c; each approved Eucharistic Prayer; CCC 1105, 1353}

GIRM calls epiclesis the Church's special invocation by which she implores the Father's power so that human gifts become Christ's Body and Blood and that the spotless Victim received in Communion may bring salvation. The definition embraces a double horizon: sanctification of the gifts and fruitful incorporation of those who receive. Later Roman prayers normally make this twofold action visible in an explicit invocation of the Holy Spirit before the institution narrative and a communion petition afterward. The Roman Canon realizes the same anaphoral necessities through a different verbal architecture.

\paragraph{The Spirit and the Son's words are not rivals.}
In Cyril's fourth-century Jerusalem catechesis, the Church asks God to send the Holy Spirit upon the gifts so that what the Spirit touches is sanctified and changed (\emph{Mystagogical Catechesis} V.7). Basil defends the Spirit's inseparable divine agency in sanctification while also appealing to the Church's received, partly unwritten liturgical practice (\emph{De Spiritu Sancto} 27, 66). Developed Eastern anaphoras often place a single extended epiclesis after the institution narrative; the postconciliar Roman compositions commonly divide explicit petition around it. These are real structural differences, not competing theologies in which one Church believes the narrative alone acts and another believes only a later invocation does.

The Catholic confession concerns the efficacy of the entire sacramental action under Christ's institution. Christ's words, pronounced by the ordained minister in his person, are consecratory; the Spirit whom the Father sends is not absent from their efficacy. The epiclesis asks for what only the triune God can do. It neither supplies human power lacking in Christ's command nor marks a second consecration.

\paragraph{How does the Roman Canon invoke?}
\latin{Quam oblationem} petitions God to order and bless the oblation so that it become Christ's Body and Blood. It is epikletic in function, yet it does not name the Holy Spirit and must not be quoted as though it used the explicit pneumatic syntax of Eucharistic Prayers II--IV\@. After the institution narrative, \latin{Supplices te rogamus} asks that participation at the earthly altar bring heavenly blessing and grace. It has a communion horizon and heavenly-altar imagery, but it is not merely a missing Eastern formula in disguise. The late-fourth-century \emph{De sacramentis} IV.5--6 securely parallels this Western sequence from the pre-institution petition through the heavenly-altar plea.

The difference teaches a methodological rule. ``Epiclesis'' identifies an anaphoral action; it does not force all prayers into one sentence pattern. An exposition should ask what is sanctified, who acts, and toward what communion the petition tends before judging a prayer by the position of a keyword.

\ordermovement{Institution Narrative and Consecration}{The Lord's Command Effects the Sacramental Memorial}{Fixed within each approved Eucharistic Prayer; showing, genuflections, bells, and incense are governed by the Missal and local use}{GIRM 79d, 150--151; \latin{Ordo Missae}; each approved Eucharistic Prayer}

Within the prayer addressed to the Father, the Priest narrates what Christ did and pronounces the Lord's words over bread and chalice. He shows each consecrated species to the people and genuflects in adoration. A bell may signal the epiclesis and each showing; incense may honor the Host and chalice. These actions manifest and adore what Christ effects. They neither cause the conversion nor interrupt the one address so that a self-contained devotional ``elevation rite'' replaces it.

\paragraph{A liturgical synthesis of the scriptural witnesses.}
Matthew~26:26--29, Mark~14:22--25, Luke~22:14--20, and 1~Corinthians~11:23--26 transmit related but non-identical accounts. Each joins meal, thanksgiving or blessing, identified gift, covenant and sacrificial self-giving, and an eschatological horizon; Luke and Paul explicitly transmit the memorial command. The Missal does not pretend to be a stenographic fifth account. It receives the Lord's actions and words through a liturgical synthesis, including the postconciliar standardization of the narrative across the Roman prayers.

Paul's context is ecclesially severe. The tradition of the Supper judges a Corinthian assembly whose meal humiliates the poor and contradicts the one body (1~Cor.~11:17--34). Proclaiming the Lord's death cannot be separated from discerning his Body and receiving worthily. Consecration is never magic performed upon isolated matter for private possession. It establishes the sacrificial banquet in which the Church must become truthful to the gift she receives.

\paragraph{Whose speech consecrates?}
Ambrose contrasts ordinary created bread with the sacrament produced by Christ's words and repeatedly asks whose speech acts: not unaided human speech, but the Lord Jesus speaking through the minister (\emph{De sacramentis} IV.4.14--17; IV.5.21--24). Aquinas gives the mature scholastic account: the Priest speaks in Christ's person, and Christ's words are the sacramental form (\emph{ST} III, q.~78, aa.~1--4; q.~82, a.~1). The minister's faithfulness, the Church's prayer, and the Spirit's action belong to the one sacramental economy; none turns the Priest into the source of divine power.

By consecration the substances of bread and wine become Christ's Body and Blood while the sensible species remain. Trent calls this unique conversion transubstantiation (Session XIII, ch.~4; canon~2), and the Catechism receives the same judgment (CCC 1376). The two consecrations sacramentally represent Christ's sacrificial self-gift under the signs of separated Body and Blood. They do not kill him again, divide the risen Christ, or require a later ritual act to reunite him. Under either species the whole Christ is present by concomitance; under both, the sign of the Eucharistic banquet is more fully manifest.

The Mass consequently neither repeats Calvary nor supplies another victim. Christ's one sacrifice is sacramentally made present and offered in an unbloody manner (Trent, Session XXII, ch.~2). The historical Last Supper, Passion, Resurrection, Ascension, and heavenly intercession remain unrepeatable events; sacramental memorial gives the Church a real participation in their saving unity.

\ordermovement{\latin{Mysterium fidei} and the Memorial Acclamation}{The Consecrated Assembly Proclaims the Paschal Mystery}{Fixed priestly invitation after the chalice showing; one approved congregational acclamation, with the children's prayers governed by their own book}{\latin{Ordo Missae}; GIRM 34, 151}

After the chalice showing and genuflection the Priest introduces the acclamation. The people then proclaim Christ's death, Resurrection, and future coming according to one of the approved forms. The acclamation is directed to Christ or proclaims him; it is not addressed to the consecrated species and should not be altered into an act of direct Eucharistic adoration. Adoration is already bodily present in the genuflections and the assembly's faith, but the received verbal action at this point is Paschal proclamation.

\paragraph{A phrase relocated, a voice newly assigned.}
In the Roman Canon before the reform, \latin{mysterium fidei} stood inside the chalice formula. Paul VI's apostolic constitution \emph{Missale Romanum} explains its removal from that position and the provision of the people's acclamation. The change should be described as a deliberate postconciliar reordering, not projected backward into apostolic practice or explained by an unsupported origin legend. The phrase now introduces the assembly's response after both consecrations.

First Corinthians~11:26 is the closest scriptural grammar: eating the bread and drinking the cup proclaims the Lord's death until he comes. The proclamation holds past, sacramental present, and promised advent together. The Resurrection prevents memory from arresting Christ at death; expectation prevents sacramental presence from being confused with the final unveiled coming. The ancient prayer \latin{Maranatha} in 1~Corinthians~16:22 and Revelation~22:20 supplies the eschatological pressure of Christian worship without being the textual source of each approved acclamation.

The alternation of voices matters. The Priest has pronounced Christ's words within the Father's address; the people respond with the Church's Paschal confession. Their acclamation does not ratify only the chalice formula---the Great Amen will ratify the whole prayer---and it does not begin the anamnesis, which the Priest immediately voices. It is a compact kerygma placed inside doxology and sacrifice.

\ordermovement{Anamnesis and Offering}{The Command to Remember Becomes Ecclesial Oblation}{Constitutive elements arranged in each Eucharistic Prayer's own grammar}{GIRM 79e--f; each approved Eucharistic Prayer}

After the institution command, the Church remembers Christ's saving Passion, Resurrection, Ascension, and, according to the prayer, his descent among the dead and future coming. In that obedient memorial she offers the Father the spotless Victim. The sequence is not first a mental recollection and then an unrelated gift. Christ commanded a memorial action in which his self-offering becomes the form and content of the Church's thanksgiving.

\paragraph{Memory before God.}
Biblical memorial is objective covenant action before it is private recall. At Passover, Israel keeps the saving act as a memorial through generations (Exod.~12:14); the rite makes each generation answerable to the deliverance it receives. In Scripture God also ``remembers'' his covenant, language that asks not for recovery from forgetfulness but for covenant fidelity manifested in saving action. At the Supper Christ gives the Church a new-covenant memorial centered upon his Body given and Blood poured out. First Corinthians~11 joins doing, eating and drinking, proclamation, self-examination, and expectation.

Justin's Sunday account joins thanksgiving, the gifts identified with Christ's flesh and blood, distribution, and care for those in need (\emph{First Apology} 65--67). Irenaeus argues against contempt for creation by pointing to bread and cup receiving Eucharistic thanksgiving and to human bodies nourished by them for resurrection (\emph{Against Heresies} IV.18.4--5; V.2.2--3). Neither witness supplies the current formulas. Together they prevent memorial from dissolving into ideas: created gifts, ecclesial thanksgiving, Christ's self-gift, moral communion, and bodily hope belong together.

\paragraph{Who offers whom?}
The Church can offer Christ because he first gives himself and commands the memorial. She does not possess him as a victim external to his own will. Nor does the assembly add a second sacrifice beside his. Augustine says that the redeemed city is offered to God as a universal sacrifice through the High Priest who offered himself and made the Church his Body (\emph{City of God} X.6, X.20). This is the deep grammar of the postconsecratory petitions: the Father receives the one Son, and those incorporated into him learn to become an offering.

\emph{Sacrosanctum Concilium} 48 therefore calls the faithful not only to offer the immaculate Victim through the Priest's hands but also to learn to offer themselves. \emph{Lumen gentium} 11 places their lives, praise, and apostolic work within Eucharistic offering. Self-offering is neither psychological intensity nor an independent human contribution. It is grace's ethical and ecclesial fruit: what lies upon the altar judges and reforms the people who say Amen.

The approved prayers articulate this unity differently. The Roman Canon surrounds the Victim with scriptural offerings and a heavenly-altar plea. Eucharistic Prayer III asks that the Spirit make the communicants an enduring offering. Eucharistic Prayer IV sets the oblation within the whole economy of creation and return. The special prayers connect it with reconciliation, pilgrimage, unity, and service. No one arrangement should be used to erase another.

\ordermovement{Communion Epiclesis and Intercessions}{The One Sacrifice Gathers a Communion across Space and Death}{Constitutive petitions and commemorations located according to each prayer; names and proper inserts are governed by the received text}{GIRM 79c, 79g; each approved Eucharistic Prayer; \emph{Paternas vices} (2013)}

The prayer asks that those who share Christ's Body and Blood be gathered in the Spirit, then displays the breadth of that communion: the Church throughout the world, pope and local bishop, ordained ministers and faithful, living needs, the dead, Mary, Joseph, apostles, martyrs, saints, and the final heavenly inheritance. The order and detail vary. Eucharistic Prayer I places substantial intercession before and after the institution narrative; later prayers generally concentrate it after anamnesis and offering. Structural location does not determine whether an intercession is essential or merely decorative.

\paragraph{The sacrament makes the body it signifies.}
First Corinthians~10:16--17 joins participation in Christ's Body and Blood to the many becoming one body because they share one bread. Ephesians~4:1--6 locates ecclesial unity in one Spirit, one body, one hope, one Lord, one faith, one baptism, and one Father. A communion epiclesis therefore asks for more than harmonious feeling. The Spirit conforms distinct persons to Christ's sacrificial charity without erasing their gifts or offices.

The \emph{Didache} uses scattered grain gathered into one loaf as an image for the Church gathered into the Kingdom (9.4). Ignatius of Antioch repeatedly joins one Eucharist, one altar, bishop, presbytery, and deacons as a concrete discipline of unity (\emph{Magnesians} 7; \emph{Philadelphians} 4). These early witnesses do not map directly onto every later intercession, but they expose a contradiction: sacramental reception cannot be invoked to sanctify faction, contempt, or refusal of ecclesial communion.

\paragraph{Why are offices and names remembered?}
Prayer for the pope and bishop places the local celebration within visible Catholic and apostolic communion. Their mention is intercessory, not a second consecratory formula or a magical test in which a verbal mistake alone decides validity. Particular names make communion historical and accountable; the larger phrases prevent any printed list from exhausting the Church. Only names and inserts the book permits belong in the Eucharistic Prayer. Private gratitude or public controversy does not authorize improvisation in a received ecclesial text.

Mary and the saints appear as creatures perfected by grace, members of Christ's Body, and companions whose promised inheritance the pilgrim Church seeks. The decree \emph{Paternas vices} of 1 May 2013 inserted Saint Joseph into Eucharistic Prayers II--IV; he had already entered the Roman Canon by John XXIII's 1962 decree. A physical 2008 Latin printing therefore lacks his name in II--IV even though the current juridical text includes it. The 2013 decree did not generally amend the Reconciliation, Various Needs, or children's prayers.

\paragraph{Can the prayer cross death?}
Second Maccabees~12:38--46 places intercession for the covenant dead within resurrection hope. Cyril describes commemoration of patriarchs, prophets, apostles, martyrs, bishops, and departed Christians during the Eucharistic sacrifice (\emph{Mystagogical Catechesis} V.8--10). Augustine records Monnica's request to be remembered at the altar and prays for her without presuming upon even a holy life (\emph{Confessions} IX.11--13). Eucharistic suffrage entrusts the dead to divine mercy through Christ's sacrifice; it neither canonizes every named person nor claims knowledge of judgment withheld from the Church.

Intercession thus reveals the social range of consecration. The gifts are changed so that a people may be changed; that people cannot be imagined apart from Catholic communion, suffering neighbors, the dead awaiting fullness, and the saints already glorified.

\ordermovement{\latin{Per ipsum} and the Great Amen}{Creation and Church Return to the Father through the Son}{Fixed priestly doxology with paten and chalice raised, ratified by the people's Amen}{GIRM 79h, 151; \latin{Ordo Missae}}

The Priest voices the Trinitarian doxology while elevating the paten with the Host. When a Deacon is present, he elevates the chalice beside the Priest; without a Deacon, the Priest elevates both paten and chalice. Concelebrants join the doxology as the Missal directs. The people answer Amen. The elevation is not a third showing or another consecration. It makes visible that the Victim, gifts, Church, and all glory return to the Father through Christ in the Spirit.

\paragraph{The grammar of return.}
Romans~11:33--36 concludes contemplation of divine mercy with the confession that all things are from, through, and for God. Ephesians~1:3--14 repeatedly orders the economy to the Father's glory through the Son and seals the inheritance in the Spirit. Hebrews~13:15--16 joins continual praise through Christ to concrete beneficence and sharing. The final doxology compresses this economy without making the Trinity a closing ornament: Christ is mediator, the Spirit is communion, and the Father is the source and end of the prayer.

Justin's second-century accounts already make the people's Amen the seal of presidential thanksgiving (\emph{First Apology} 65, 67). In 1~Corinthians~14:16 Paul assumes that another can answer Amen to thanksgiving only if the prayer is intelligible. The liturgical answer is therefore neither applause nor a cue that ``the important part'' has ended. It appropriates an address the people did not voice paragraph by paragraph. The more fully the prayer has named sacrifice, Church, living, dead, creation, and Kingdom, the more comprehensive the Amen must become.

\paragraph{Where has the prayer gone?}
The Preface lifted hearts and named divine works; the Sanctus placed earth within heavenly praise; the Father was implored to sanctify gifts and people; Christ's words effected sacramental presence; the assembly proclaimed his Pasch; the Church remembered and offered; the Spirit's communion widened into intercession; and the whole now returns through the Son. None of these is a detachable devotion. The Amen belongs to the entire arc.

\Needspace{10\baselineskip}
\begin{massbox}{One Eucharistic Prayer}
Thanksgiving does not merely introduce consecration; consecration does not suspend praise; memorial is already offering; offering tends toward Communion; Communion opens into the whole Church; intercession reaches the dead and the Kingdom; and the doxology returns every received gift to the Father. The people's Amen is the audible assent of the body joined to Christ's act.
\end{massbox}
